%% Method subsections for RQ4 and RQ5 (Phase 7 and 8 protocols).
\subsection{Remedy and sensitivity}\label{sec:rq4method}

RQ4 reuses the grid of RQ2 and adds no new main model. H4 is judged on the common population with the schedule ground truth and the conversation-grouped split. Beside it we report, per class, the change in the exporter share from the flow unit to the conversation unit, with a paired bootstrap interval, and call it a fall or a rise when the interval excludes zero. A class has a large residual when its exporter share at the conversation unit is at least \setnum{0.25} and the arm means of its F1 spread by at least \setnum{0.05}. For each such class we break the errors down by arm: false negatives and false positives, whether a false positive's conversation has several records, and whether exactly one of them was flagged. Leave-one-arm-out shares, without intervals, show which arm carries a residual.

The sensitivity variants change one choice at a time: the gap $T$ (\setnum{5}, \setnum{30}, \setnum{60} and \setnum{120} seconds, with conversations rebuilt by the same code, which at \SetT{} seconds reproduces the frozen tables exactly), the aggregation rule (majority, byte-weighted), the unit (attack episodes in place of conversations), the treatment of attempted attacks (folded into benign) and the labels (each release's own in place of the schedule). The last two are run twice, once scoring the primary models against the changed labels and once retraining on them, the retraining on the grouped split and on \cicids{} only. For every variant we state which of five conclusions change: H4, H2 at the flow unit, H2 at the conversation unit, the per-class verdicts of exporter against model at the conversation unit, and the per-class direction of the exporter share.

\subsection{Re-examination of published claims}\label{sec:claims}

The claims for RQ5 were fixed in a table written before any RQ5 number was computed, so that they could not be chosen after seeing which survive. It lists \docnum{11} candidate claims from the papers of Section~\ref{sec:related}, of which five can be examined with released data (Table~\ref{tab:rq5}). The others need a re-extraction of the packet captures over a timeout grid \cite{janatiidrissi2025}, flows that have not been released \cite{pinto2025,xin2026}, or benchmarks outside our scope \cite{vitorino2025}. For BCCC-CIC-IDS2017 we have not read the claim text, and NF-UNSW-NB15-v3 makes no performance comparison \cite{luay2026}.

Each examinable contrast goes through three steps. First, the comparison is rebuilt per flow the way the authors made it, on their releases with their own labels, preprocessing, features and model, using the authors' code where it is public and a line-by-line port of it where needed, refitted with seeds \SetSeeds{}. We do not claim to reproduce their exact experiment, and we state where our setup departs from theirs. Second, the per-flow gap $G$ is split into a part from classes that only one release has and, on the shared classes, a class weight part $C$ and a within-class part $W$. $C$ is the symmetric two-fold reweighting of the test records of each release to the class shares of the other, in the manner of an Oaxaca decomposition, and $W = G_{\mathrm{shared}} - C$. Third, the same contrast is computed on the conversations of the subsample of Section~\ref{sec:protocol} that have records in both releases, with the schedule ground truth, the conversation-grouped split, the claim's own model and features, and the ``any'' rule. Intervals are paired cluster bootstrap intervals over conversations with \SetBoot{} resamples. The verdict is read from the conversation-level gap in the claim's own metric. A claim survives when the gap has the published sign and its interval excludes zero, reverses when the interval lies entirely on the other side, and is not supported otherwise. A claim of no difference survives when the interval lies within \setnum{0.01} of zero and is contradicted when it excludes zero. A class with fewer than \SetMinTest{} test conversations in any seed is not estimable.
